\chapter*{Tóm tắt}
%\markboth{TÓM TẮT}
Đề tài “Cải thiện khả năng tự phục hồi cho hệ thống Kubernetes sử dụng học tăng cường” tập trung nghiên cứu và xây dựng mô hình ứng dụng Reinforcement Learning nhằm tối ưu quá trình phát hiện và xử lý sự cố trong Kubernetes cluster.

Trong mô hình đề xuất, tác nhân học tăng cường (RL Agent) sẽ liên tục quan sát trạng thái hệ thống như mức sử dụng CPU, RAM, trạng thái pod, node và các chỉ số hiệu năng khác để đưa ra hành động phù hợp như restart pod, migrate workload hoặc phân phối lại tài nguyên. Sau mỗi hành động, hệ thống sẽ phản hồi thông qua reward nhằm giúp agent học được chiến lược phục hồi tối ưu theo thời gian.

Đề tài hướng đến mục tiêu:
\begin{itemize}
    \item Giảm thời gian gián đoạn dịch vụ khi xảy ra lỗi.
    \item Nâng cao khả năng tự động thích nghi của hệ thống.
    \item Tối ưu hiệu quả sử dụng tài nguyên trong cluster.
    \item Cải thiện tính ổn định và độ sẵn sàng của ứng dụng.
\end{itemize}
Thông qua việc tích hợp Reinforcement Learning vào cơ chế self-healing của Kubernetes, hệ thống có thể phản ứng linh hoạt hơn trước các sự cố phức tạp thay vì chỉ phụ thuộc vào các luật xử lý cố định truyền thống.